% Options for packages loaded elsewhere
\PassOptionsToPackage{unicode}{hyperref}
\PassOptionsToPackage{hyphens}{url}
\documentclass[
]{article}
\usepackage{xcolor}
\usepackage[margin=1in]{geometry}
\usepackage{amsmath,amssymb}
\setcounter{secnumdepth}{-\maxdimen} % remove section numbering
\usepackage{iftex}
\ifPDFTeX
  \usepackage[T1]{fontenc}
  \usepackage[utf8]{inputenc}
  \usepackage{textcomp} % provide euro and other symbols
\else % if luatex or xetex
  \usepackage{unicode-math} % this also loads fontspec
  \defaultfontfeatures{Scale=MatchLowercase}
  \defaultfontfeatures[\rmfamily]{Ligatures=TeX,Scale=1}
\fi
\usepackage{lmodern}
\ifPDFTeX\else
  % xetex/luatex font selection
\fi
% Use upquote if available, for straight quotes in verbatim environments
\IfFileExists{upquote.sty}{\usepackage{upquote}}{}
\IfFileExists{microtype.sty}{% use microtype if available
  \usepackage[]{microtype}
  \UseMicrotypeSet[protrusion]{basicmath} % disable protrusion for tt fonts
}{}
\makeatletter
\@ifundefined{KOMAClassName}{% if non-KOMA class
  \IfFileExists{parskip.sty}{%
    \usepackage{parskip}
  }{% else
    \setlength{\parindent}{0pt}
    \setlength{\parskip}{6pt plus 2pt minus 1pt}}
}{% if KOMA class
  \KOMAoptions{parskip=half}}
\makeatother
\usepackage{graphicx}
\makeatletter
\newsavebox\pandoc@box
\newcommand*\pandocbounded[1]{% scales image to fit in text height/width
  \sbox\pandoc@box{#1}%
  \Gscale@div\@tempa{\textheight}{\dimexpr\ht\pandoc@box+\dp\pandoc@box\relax}%
  \Gscale@div\@tempb{\linewidth}{\wd\pandoc@box}%
  \ifdim\@tempb\p@<\@tempa\p@\let\@tempa\@tempb\fi% select the smaller of both
  \ifdim\@tempa\p@<\p@\scalebox{\@tempa}{\usebox\pandoc@box}%
  \else\usebox{\pandoc@box}%
  \fi%
}
% Set default figure placement to htbp
\def\fps@figure{htbp}
\makeatother
\setlength{\emergencystretch}{3em} % prevent overfull lines
\providecommand{\tightlist}{%
  \setlength{\itemsep}{0pt}\setlength{\parskip}{0pt}}
\usepackage{bookmark}
\IfFileExists{xurl.sty}{\usepackage{xurl}}{} % add URL line breaks if available
\urlstyle{same}
\hypersetup{
  pdftitle={Relazione Progetto Parallelo},
  pdfauthor={Sebastiano Patton},
  hidelinks,
  pdfcreator={LaTeX via pandoc}}

\title{Relazione Progetto Parallelo}
\author{Sebastiano Patton}
\date{2026-04-20}

\begin{document}
\maketitle

{
\setcounter{tocdepth}{2}
\tableofcontents
}
\section{Introduzione}\label{introduzione}

Lo scopo del progetto è progettare e implementare un software c++ e
OpenMp parallelo per poi fare un'analisi delle prestazioni analizzando
l'impatto della variazione dei parametri. Ho deciso di scrivere un
programma che calcola l'insieme di Mandelbrot, ovvero un insieme di
numeri complessi c per i quali la successione definita da:
\[\begin{equation}\begin{cases}Z_0 = 0 \\Z_{n+1} = Z_n^2 + c\end{cases}\end{equation}\]
è limitata. Questo programma va quindi a calcolare il numero di
iterazioni necessarie per far tendere all'infinito il punto, ciò avviene
quando la distanza dal centro del punto z è maggiore di 2 ovvero quando:
\[|Z|^2 = Z_x^2 +Z_y^2 > 4\]
\pandocbounded{\includegraphics[keepaspectratio]{/Progetto_Parallelo/Relazione/grafici/compute.png}}

quando il codice raggiunge il numero massimo di iterazioni per un punto
si considera tale punto parte dell'insieme e in ogni caso va a salvare
il numero di iterazioni svolte per ogni pixel dell'immagine. Nel codice
si utilizza una variabile chiamata tmp per conservare il vecchio valore
di Zx in quanto lo ne ho bisogno per calcolare la nuova Zy. Per generare
l'immagine in output il codice calcola il passo lungo X e lungo Y,
ovvero va a calcolare a quanto spazio del piano immaginario corrisponde
un pixel e salva questi valori come due costanti:
\[\frac{X_{max} - X_{min}}{X_{p}}\] \[\frac{Y_{max} - Y_{min}}{Y_{p}}\]
Con Xp e Yp che indicano il numero di pixel lungo X e lungo Y con cui
creo l'immagine. Il codice userà i passi dx e dy per calcolare le
coordinate matematiche di ogni pixel ovvero cx e cy che andrà poi ad
usare per il calcolo del numero di iterazioni. Per farlo somma al valore
minimo del piano immaginario che consideriamo il numero della riga o
della colonna moltiplicato per il passo rispettivo (dy per le righe e dx
per le colonne). Tramite la libreria grafica il numero viene
interpretato e gli si assegna un colore.

\pandocbounded{\includegraphics[keepaspectratio]{/Progetto_Parallelo/Relazione/grafici/render.png}}

Il codice creerà nella cartella output delle immagini come questa:
\pandocbounded{\includegraphics[keepaspectratio]{/Progetto_Parallelo/Relazione/grafici/test_par_4096.png}}

Il motivo per cui ho scelto l'insieme di mandelbrot è che si tratta di
un problema ``embarassingly parallel'', questo perché il calcolo del
valore corrispondente ad ogni pixel è completamente indipendente dagli
altri rendendolo un codice facilmente parallelizzabile. Questo lo rende
un programma adatto per analizzare le diverse configurazioni di
parallelizzazione come il variare del numero dei pixel, lo scheduling e
il numero di thread.

Per fare questi test ho utilizzato il mio computer che ha un processore
amd ryzen 3600xt con 6 core e 12 processori logici totali insieme a 16gb
di ram. Per fare l'analisi delle prestazioni il dato più importante di
cui necessito è il tempo impiegato ad eseguire il programma, che non
deve includere il tempo di scrittura e lettura dai file contenenti i
parametri necessari per il funzionamento del codice. Per raccogliere
questo dato ho usato il comando omp\_get\_wtime() che viene chiamato
subito prima e subito dopo l'esecuzione del codice andando così a
salvare gli orari in cui è stato chiamato per poi calcolare il tempo con
una sottrazione (tempo finale - tempo inizale)

\section{Parallelizzazione}\label{parallelizzazione}

Nel mio codice sono presenti due cicli for annidati, il primo (il più
esterno) scorre le righe mentre il secondo scorre le colonne. Per
parallelizzare il codice in modo efficiente faccio svolgere in parallelo
il ciclo più esterno dato che: se il codice svolgesse in parallelo il
ciclo interno si andrebbero a creare e distruggere tantissimi thread
andando così a generare un overhead che renderebbe il codice parallelo
più lento del codice seriale. Un altro problema che occorrerebbe sarebbe
che i vari thread userebbero gli stessi pixel in quanto non sarebbero
dichiarate bene quali variabili sono condivise e quali sono private.
Parallelizzando il ciclo più esterno si impostano come variabili
condivise: numero di iterazioni, dimensione in pixel dell'immagine, le
dimensioni della sezione di piano immaginario che voglio analizzare e i
passi che ho calcolato (dx e dy), e come variabili private ho impostato:
``row'' ovvero l'indice del ciclo for principale (righe), ``col''
l'indice del ciclo annidato (colonne) ed infine le coordinate cx e cy
correnti. Per il tipo di codice che ho scelto è quindi impossibile
andare a cambiare quali variabili sono private e quali sono condivise,
perché andrebbe a creare dei conflitti che non farebbero risultare
correttamente l'immagine.

\section{Analisi delle prestazioni}\label{analisi-delle-prestazioni}

\subsection{Confronto seriale e
parallelo}\label{confronto-seriale-e-parallelo}

Come primo test ho deciso di fare lo studio delle performance al variare
della dimensione dell'immagine, in serie e in parallelo (scheduling:
dinamico); e ho osservato il variare delle prestazioni con l'aumento del
numero di pixel da calcolare. Si nota che per risoluzioni basse ovvero
512x512 e 1024x1024 il tempo diesecuzione è molto simile per seriale e
parallelo ma, salendo con la qualità delle immagini, si vede aumentare
notevolmente il tempo del seriale (arrivando fino a 10.6 secondi per
4096x4096 e 41.5 secondi per 8192x8192) mentre il parallelo mantiene un
esecuzione molto veloce impiegandoci addirittura il 90\% del tempo in
meno.

\pandocbounded{\includegraphics[keepaspectratio]{/Progetto_Parallelo/Relazione/grafici/serievsparallelo.png}}

\subsection{Studio Performance al variare del numero di
thread}\label{studio-performance-al-variare-del-numero-di-thread}

I thread sono l'unità minima di esecuzione che può essere gestita da un
sistema operativo e il multithreading è il fondamento della
parallelizzazione, questo perchè i vari core della cpu per svolgere un
programma in parallelo creano vari thread su cui svolgono operazioni
diverse in contemporanea permettendo quindi elevate velocità di calcolo.
In base al numero di thread che impongo alla cpu di creare si può
individuare quale sia la quantità ideale per il funzionamento del
codice, questo perchè se ho troppi pochi thread la cpu non andrà al
massimo, mentre se ne ho troppi andrei a continuare a creare thread
nuovi andando a rallentare l'esecuzione del programma. Ho notato che ho
il funzionamento migliore quando il numero di thread è 12 ovvero pari al
numero di processori logici della mia cpu. Ho fatto i test con
1,12,512,2048,8192 e 16384 thread e una risoluzione di 4096 e 512,In
entrambi i casi il più lento è chiaramente stato quello con 1 thread in
quanto è l'equivalente di svolgere il codice in seriale, mentre ho
notato un aumento significativo dei tempi di esecuzione da 8192 thread
in poi.

\pandocbounded{\includegraphics[keepaspectratio]{/Progetto_Parallelo/Relazione/grafici/numerodithread.png}}

Il grafico ha una scala logaritmica per l'asse delle ascisse

\subsubsection{Speedup e Overhead}\label{speedup-e-overhead}

Con lo studio delle performanca al variare del numero dei thread ho
potuto analizzare anche due caratteristiche importanti per la
parallelizzazione ovvero Speedup e Overhead:

-Speedup: Di quante volte aumenta la velocità di esecuzione
all'aumentare del numero di thread
\(Speedup = \frac{t_{seriale}}{t_{nthread}}\)

\pandocbounded{\includegraphics[keepaspectratio]{/Progetto_Parallelo/Relazione/grafici/speedup.png}}

Si nota che lo speedup aumenta fino ad avere circa 12 thread,
successivamente lo speedup diminuisce sempre di più. Il codice rimane
comunque più veloce in parallelo. Dato che l'immagine 512x512 è molto
piccola il codice suddivide il carico in task di dimensioni ridotte ad
ogni thread andando andando così ad aumentare notevolmente l'overhead e
diminuendo più velocemente lo speedup.

-Overhead: è il tempo perso ad assegnare i valori richiesti ad ogni
thread \(Overhead = t_{nthread} - \frac{t_{1thread}}{n_{thread}}\)

\pandocbounded{\includegraphics[keepaspectratio]{/Progetto_Parallelo/Relazione/grafici/overhead.png}}

Dal grafico dell'overhead è chiaro che un eccessivo aumento del numero
di thread va solo a causare dei rallentamenti nell'esecuzione del codice
in quanto aumenta la quantità di richieste di dati andando a rallentare
l'esecuzione.

\subsection{Studio delle performance in base allo
scheduling}\label{studio-delle-performance-in-base-allo-scheduling}

Lo scheduling è il modo in cui le varie iterazioni dei cicli for vengono
forniti ai thread. Ci sono tre tipi di scheduling: static, dynamic e
guided.

-Static: divide le iterazioni del ciclo for in blocchi di uguali
dimensioni e liassegna ognuno ad un thread. Quando un thread finisce di
calcolare un blocco, gli si fornisce un altro blocco di uguali
dimensioni. Solitamente le dimensioni di questo blocco sono N/threads.
Questo tipo di scheduling non fa caso al tempo impiegato da ciascun
thread per risolvere il blocco, questo può causare sbilanciamenti nel
carico in quanto alcuni blocchi sono più pesanti di altri e di
conseguenza richiedono più tempo per essere calcolati. Un vantaggio è
che ha il minor overhead di tutti in quanto tende a dividere il numero
di iterazioni per il numero di thread e assegnare uno di questi blocchi
ad ognuno di questi thread, così facendo non sono richieste
comunicazioni e quindi non c'è overhead.

-Dynamic: L'assegnazione delle iterazioni avviene durante l'esecuzione
del ciclo, quindi si assegnano dei blocchi ad ogni thread e appena uno
dei thread finisce di calcolare un blocco ne riceve un altro da
calcolare. La dimensione di default dei chuank è 1. Così facendo non ci
sono blocchi inattivi ed è l'ideale per codici con blocchi con tempi
molto variabili come nel caso del mio codice. Data la quantità di
richieste di blocchi da assegnare lo scheduling dynamic presenta
abbastanza overhead.

-Guided: Anche qua l'assegnazione delle iterazioni avviene
durantel'esecuzione del ciclo solo che i blocchi che vengono assegnati
ai vari thread diminuiscono di dimensioni mano a mano che proseguono le
assegnazioni. I blocchi assegnati sono di dimensioni proporzionali al
numero di operazioni rimantenti diviso il numero di thread. Lo scopo di
questo scheduling è diminuire l'overhead iniziale diminuendo le
richieste all'inizio e aumentandole man mano durante l'esecuzione.

\pandocbounded{\includegraphics[keepaspectratio]{/Progetto_Parallelo/Relazione/grafici/scheduling.png}}

Ho fatto i test con risoluzioni di 512x512, 1024x1024, 2048x2048,
4096x4096 e 8192x8192. Le risoluzioni più alte mettono in evidenza le
differenze di tempo richieste per il calcolo delle iterazioni in base
allo scheduling e chiaramente: lo scheduling dynamic è il più veloce, il
guided è una via di mezzo mentre lo static è il più lento. Tale
risultato è corente con l'analisi visiva dell'immagine per generata per
capirlo, le sezioni blu sono quelle che richiedono meno iterazioni per
tendere all'infinito mentre quelle nere sono quelle che ne richiedono di
più. le sezioni nere sono tutte abbastanza vicine tra di loro e quindi
per uno scheduling static o guided è probabile che nei grandi pacchetti
(iniziali nel caso del guided) che vanno a calcolare ce ne siano alcuni
con la maggior parte di linea che richiedono poche iterazioni mentre
altre che ne richiedono molte in quanto hanno righe principalmente nere.
Questo rende lo scheduling dynamic il più veloce dato che fin da subito
prende pacchetti piccoli e nonostante il maggior overhead riesce
comunque ad essere più veloce.

\subsubsection{Variazione del numero dei
chunk}\label{variazione-del-numero-dei-chunk}

I chunk sono la dimensione dei pacchetti che vengono assegnati ai thread
ovvero la quantità di righe che assegno ad ogni thread, variandole vado
a modificare la velocità dei vari scheduling. Guardando il grafico si
nota che nonostante l'aumento dell'overhead, tenere le dimensioni dei
chunk basse aiuta a diminuire i tempi di esecuzione per scheduling
dynamic e static mentre lo scheduling guided rimane abbastanza costante
fino a chunk da 128 righe. questo accade perchè nel caso dello
scheduling guided la dimensione dei chunk che si imposta indica il
valore minimo che possono ragigungere i chunk con l'andare del
programma, per questo se il valore minimo è minore di quanto gli avrebbe
assegnato in ogni caso il numero di chunk non influisce.ù

\pandocbounded{\includegraphics[keepaspectratio]{/Progetto_Parallelo/Relazione/grafici/chunk.png}}

\section{Conclusioni}\label{conclusioni}

Visti i risultati dei vari test ritengo che per il miglior metodo per
parallelizzare l'insieme di mandelbrot su un computer con 6 core sia di
utilizzare un scheduling dinamico (con dimensioni dei chunk di 1) e 12
thread per parallelizzare il ciclo for più esterno, ovvero quello delle
righe. Queste impostazioni permettono di abbattere i tempi di esecuzione
del 90\% per immagini molto pesanti rendendolo quindi estremamente più
efficiente.

\end{document}
